\bookchapter{Gradient descent}\label{ch:gd}

\section*{Explanation}
We cannot always solve $f'(x)=0$ by hand, especially with millions of variables. Instead we walk downhill step by step. Start at a guess $x_0$ and repeat
\[
x_{n+1}=x_n-\eta\,f'(x_n),
\]
where $\eta>0$ is the \emph{learning rate} (step size). For several variables $\mathbf w=(w_1,\dots,w_k)$:
\[
\mathbf w_{\text{new}}=\mathbf w_{\text{old}}-\eta\,\nabla L(\mathbf w_{\text{old}}).
\]
The minus sign is crucial: the gradient points uphill, so we go the opposite way.

\begin{example}
Minimise $f(x)=x^2$ with $f'(x)=2x$, start $x_0=5$, $\eta=0.1$:
\begin{align*}
x_1&=5-0.1\cdot10=4,\\
x_2&=4-0.1\cdot 8=3.2,\\
x_3&=3.2-0.1\cdot6.4=2.56.
\end{align*}
Each step multiplies $x$ by $0.8$, so $x_n\to0$, which is the minimum.
\end{example}

\begin{example}[Learning rate too large]
With $\eta=1.1$ on $f(x)=x^2$: $x_1=5-1.1\cdot10=-6$, $x_2=-6+13.2=7.2$, \dots Each step multiplies $x$ by $-1.2$, so the iterates jump back and forth and grow. The method \emph{diverges}.
\end{example}

\begin{example}[Two variables]
$f(x,y)=x^2+y^2$, $\nabla f=(2x,2y)$, start $(3,4)$, $\eta=0.25$:
\[
(3,4)-0.25\cdot(6,8)=(1.5,\,2).
\]
\end{example}

\begin{aibox}
Training a neural network is gradient descent on the loss, with the weights as variables. Too small a learning rate means slow training, too large means unstable training.
\end{aibox}

\section*{Practice problems}
\begin{enumerate}[label=\textbf{\thechapter.\arabic*}]
  \item Let $f(x)=(x-4)^2$ and start at $x_0=0$ with $\eta=0.25$. Compute $x_1,x_2,x_3$. Where are the iterates heading?
  \item Repeat 8.1 with $\eta=1$. What happens?
  \item Let $f(x,y)=x^2+2y^2$. Start at $(2,1)$ with $\eta=0.1$ and compute one step of gradient descent.
  \item Let $L(w)=(2w-6)^2$, start at $w_0=0$, $\eta=0.05$. Compute $w_1$. What is the true minimiser?
  \item Explain in your own words what happens if the learning rate is (a) very small, (b) very large.
  \item[\textbf{8.6}] \texttt{[code]} Write a Python loop that performs $20$ steps of gradient descent for $f(x)=(x-4)^2$ starting at $x_0=0$ and prints $x$ after each step. Try different values of $\eta$.
  \item[\textbf{8.7}$\star$] For $f(x)=x^2$ show that gradient descent gives $x_{n+1}=(1-2\eta)\,x_n$. For which values of $\eta$ does $x_n\to0$?
\end{enumerate}

% =====================================================
